\subsection{DIRECT SUM OF ROOT SYSTEMS}

Let \(V\) be a vector space over \(k\) that is the direct sum of a family \((V_i)_{1 \leqslant i \leqslant r}\) of vector spaces. Identify \(V^*\) with the direct sum of the \(V_i^*\) . For all \(i\) , let \(R_i\) be a root system in \(V_i\) . Then \(R = \bigcup_i R_i\) is a root system in \(V\) whose inverse system is \(R^- = \bigcup_i R_i^-\) ; the canonical bijection from \(R\) to \(R^-\) extends, for all \(i\) , the canonical bijection from \(R_i\) to \(R_i^-\) . The set \(R\) is called the direct sum of the root systems \(R_i\) . Let \(\alpha \in R_i\) . If \(j \neq i\) , the kernel of \(\alpha^-\) contains \(V_j\) , so \(s_\alpha\) induces the identity on \(V_j\) ; on the other hand, \(k\alpha \subseteq V_i\) , so \(s_\alpha\) leaves \(V_i\) stable. These remarks show that \(W(R)\) can be identified with \(\prod_{i=1}^{r} W(R_i)\) .

A root system R is said to be irreducible if \(R \neq \varnothing\) and if R is not the direct sum of two non-empty root systems.

PROPOSITION 5. Let V be a vector space over k that is the direct sum of vector spaces \(V_{1}, \ldots, V_{r}\) . Let R be a root system in V. Put \(R_{i} = R \cap V_{i}\) . The following three conditions are equivalent:

\begin{enumerate}
\def\labelenumi{(\roman{enumi})}
\item
  the \(V_{i}\) are stable under W(R);
\item
  \(R \subseteq V_{1} \cup V_{2} \cup \cdots \cup V_{r};\)
\item
  for all \(i\) , \(\mathbb{R}_i\) is a root system in \(\mathrm{V}_i\) , and \(\mathbb{R}\) is the direct sum of the \(\mathbb{R}_i\) .
\item
  \(\Longrightarrow\) (i): this follows from what we said at the beginning of this number.
\item
  \(\Longrightarrow\) (ii): assume that the \(V_{i}\) are stable under \(W(R)\) . Let \(\alpha \in R\) and let H be the kernel of \(\alpha\) . By Prop. 3 of Chap. V, § 2, no. 2, each \(V_{i}\) is the sum of a subspace of H and a subspace of \(k\alpha\) . Hence one of the \(V_{i}\) contains \(k\alpha\) , so \(\alpha \in V_{1} \cup V_{2} \cup \cdots \cup V_{r}\) .
\item
  \(\Longrightarrow\) (iii): if condition (ii) is satisfied, \(\mathbf{R}_i\) generates \(\mathbf{V}_i\) for all \(i\) , so \(\mathbf{R}_i\) is a root system in \(\mathbf{V}_i\) (Prop. 4). It is clear that \(\mathbf{R}\) is the direct sum of the \(\mathbf{R}_i\) .
\end{enumerate}

COROLLARY. Let R be a root system in V. The following conditions are equivalent:

\begin{enumerate}
\def\labelenumi{(\roman{enumi})}
\item
  R is irreducible:
\item
  the W(R)-module V is simple;
\item
  the W(R)-module V is absolutely simple.
\item
  \(\Longleftrightarrow\) (i): this follows from Prop. 5 and Maschke's Theorem (Chap. V, Appendix, Prop. 2).
\item
  \(\Longleftrightarrow\) (ii): this follows from Prop. 1 of Chap. V, § 2, no. 1.
\end{enumerate}

PROPOSITION 6. Every root system R in V is the direct sum of a family \((\mathrm{R}_{i})_{i\in I}\) of irreducible root systems that is unique up to a bijection of the index set.

The existence of the \(R_{i}\) is proved by induction on Card R: if R is non-empty and not irreducible, R is the direct sum of two root systems \(R'\) , \(R''\) such that \(Card R' < Card R\) , \(Card R'' < Card R\) , and the induction hypothesis applies to \(R'\) and \(R''\) . To prove uniqueness, it suffices to prove that, if R is the direct sum of \(R'\) and \(R''\) , every \(R_i\) is necessarily contained either in \(R'\) or in \(R''\) . Let \(V', V'', V'_i, V''_i\) be the vector subspaces of V generated by \(R', R'', R' \cap R_i, R'' \cap R_i\) . Since the sum \(V' + V''\) is direct, the sum \(V'_i + V''_i\) is direct. Since \(R_i \subseteq R' \cup R''\) , \(R_i\) is the direct sum of the root systems \(R' \cap R_i\) and \(R'' \cap R_i\) ; hence either \(R' \cap R_i = \varnothing\) or \(R'' \cap R_i = \varnothing\) , which proves the assertion.

The \(R_{i}\) are called the irreducible components of R. For any non-zero scalars \(\lambda_{i}\) , the union of the \(\lambda_{i}R_{i}\) is a root system in V, whose inverse system is the union of the \(\lambda_{i}^{-1}R_{i}\) , and whose Weyl group is \(\mathrm{W}(\mathrm{R})\) .

PROPOSITION 7. Let R be a root system in V, (R \(_{i}\) ) the family of its irreducible components, V \(_{i}\) the vector subspace of V generated by R \(_{i}\) , B the invariant symmetric bilinear form on V defined in Prop. 3, and B' a symmetric bilinear form on V invariant under W(R). Then the V \(_{i}\) are pairwise orthogonal with respect to B', and, for all i, the restrictions of B and B' to V \(_{i}\) are proportional.

If \(v_{i} \in \mathrm{V}_{i}\) , \(v_{j} \in \mathrm{V}_{j}\) , \(i \neq j\) , and if \(w \in \mathrm{W}(\mathrm{R}_j)\) , then

\[
\mathrm{B} ^ {\prime} (v _ {i}, w (v _ {j})) = \mathrm{B} ^ {\prime} (v _ {i}, v _ {j}),
\]

which shows that \(w(v_{j}) - v_{j}\) is orthogonal to \(v_{i}\) with respect to \(B'\) . Since \(V_{j}\) is irreducible for \(\mathrm{W}(\mathrm{R}_{j})\) , it is generated by the \(w(v_{j}) - v_{j}\) , and it is therefore orthogonal to \(V_{i}\) .

The fact that the restrictions of B and B' to each of the \(V_{i}\) are proportional follows from Prop. 1 of Chap. V, § 2, no. 1.

Remark. Choose a scalar product on \(V_{R}\) invariant under \(W(R)\) . It is then possible to speak of the length of a root and the angle between two roots. Prop. 7 shows that this angle is independent of the choice of scalar product, as is the ratio of the lengths of two roots, provided they belong to the same irreducible component of R.

